\chapter{The multiplication table (Shimura 3.24)}\label{chap:multiplication-table}

Because the $\mathrm{GL}_2$ Hecke algebra $\mathcal{H}_2$ is commutative
(Theorem~\ref{thm:commring-hecke-algebra}), its structure is completely pinned down by how the
generators $T(a,d)$ multiply. This chapter formalises Shimura's Theorem~3.24 [Sh, Thm~3.24]: the
identity $T(1)=1$, the prime-power telescoping relation $T(1,p^k)=T(p^k)-T(p,p)\,T(p^{k-2})$, the
recurrence and product formula for the prime-power operators $T(p^r)$, full multiplicativity
$T(m)\,T(n)=\sum_{d\mid(m,n)} d\,T(d,d)\,T(mn/d^2)$, and the degree formulas
$\deg T(p^k)=1+p+\dots+p^k$ and $\deg T(m)=\sigma_1(m)$. The arguments are arithmetic and
combinatorial: telescoping sums of diagonal classes, the central role of the scalar class
$T(p,p)$, coset-counting for the structure constants, and the Chinese Remainder Theorem for the
coprime factorisation.

\begin{lemma}[$T(1)=1$]\label{lem:T-sum-one}
\lean{HeckeRing.GL2.T_sum_one}
\uses{def:T-sum, def:T-ad}
\leanok
In $\mathcal{H}_2$ one has $T(1)=1$, the multiplicative identity of the Hecke algebra.
\end{lemma}

\begin{proof}
\uses{def:T-sum, def:T-ad, lem:T-ad-one-one}
\leanok
The only positive divisor of $1$ is $1$ itself, so the defining sum
$T(1)=\sum_{a\mid 1} T(a,1/a)$ collapses to the single term $T(1,1)$. By
Lemma~\ref{lem:T-ad-one-one} this equals the identity of $\mathcal{H}_2$.
\end{proof}

\begin{theorem}[Shimura 3.24(2): telescoping relation]\label{thm:T-ad-one-ppow-eq}
\lean{HeckeRing.GL2.T_ad_one_ppow_eq}
\uses{def:T-sum, def:T-ad, def:T-pp}
\leanok
Let $p$ be a prime and $k\ge 2$. Then in $\mathcal{H}_2$
\[
  T(1,p^k) \;=\; T(p^k) \;-\; T(p,p)\,T(p^{k-2}).
\]
\end{theorem}

\begin{proof}
\uses{def:T-sum, def:T-ad, def:T-pp}
\leanok
Expand $T(p^k)=\sum_{i=0}^{\lfloor k/2\rfloor} T(p^i,p^{k-i})$ over the diagonal classes whose
elementary divisors are powers of $p$. Multiplying the scalar class $T(p,p)$ into a diagonal
class shifts both elementary divisors up by one, $T(p,p)\,T(p^j,p^{d})=T(p^{j+1},p^{d+1})$, so
$T(p,p)\,T(p^{k-2})=\sum_{j} T(p^{j+1},p^{k-1-j})$, which is exactly the expansion of $T(p^k)$
with the leading $i=0$ term $T(1,p^k)$ removed. Subtracting the two expansions telescopes to the
single surviving term $T(1,p^k)$, giving the stated identity. (Equivalently, the displayed sum is
proved by reindexing the shifted sum via $i=j+1$ and peeling off the bottom term.)
\end{proof}

\begin{theorem}[Key prime computation]\label{thm:T-sum-prime-mul-T-ad}
\lean{HeckeRing.GL2.T_sum_prime_mul_T_ad}
\uses{lem:T-sum-prime, def:T-ad, def:T-pp, def:hecke-multiplicity}
\leanok
Let $p$ be a prime and $k\ge 1$. Then
\[
  T(p)\,T(1,p^k) \;=\; T(1,p^{k+1}) \;+\; c\,T(p,p^k),
  \qquad
  c \;=\;
  \begin{cases}
    p+1, & k=1,\\
    p,   & k\ge 2.
  \end{cases}
\]
\end{theorem}

\begin{proof}
\uses{lem:T-sum-prime, def:T-ad, def:T-pp, def:hecke-multiplicity}
\leanok
Since $T(p)=T(1,p)$ (Lemma~\ref{lem:T-sum-prime}), the product is the convolution of the basis
elements attached to $\operatorname{diag}(1,p)$ and $\operatorname{diag}(1,p^k)$. By the product
formula for the Hecke ring, its coefficient at a double coset $A$ is the structure constant
$\mu$ counting pairs of left-coset representatives whose product lands in a left coset of $A$.
A determinant computation forces every such $A$ to have total elementary divisor product
$p^{k+1}$, and a divisibility analysis of the first elementary divisor (using that an integral
matrix of determinant $1$ has coprime first column) shows that the only double cosets that can
occur are $\operatorname{diag}(1,p^{k+1})$ and $\operatorname{diag}(p,p^k)$. Counting the
overlapping left cosets gives multiplicity $1$ at $\operatorname{diag}(1,p^{k+1})$ and
multiplicity $c$ at $\operatorname{diag}(p,p^k)$, where the count splits according to whether
$k=1$ (the two contributing cosets coincide, giving $p+1$) or $k\ge 2$ (giving $p$). All other
coefficients vanish, which is the claimed identity.
\end{proof}

\begin{theorem}[Prime-power recurrence]\label{thm:T-sum-ppow-recurrence}
\lean{HeckeRing.GL2.T_sum_ppow_recurrence}
\uses{thm:T-ad-one-ppow-eq, thm:T-sum-prime-mul-T-ad, lem:T-sum-prime, def:T-pp}
\leanok
Let $p$ be a prime and $k\ge 1$. Then
\[
  T(p^{k+1}) \;=\; T(p)\,T(p^k) \;-\; p\,\bigl(T(p,p)\,T(p^{k-1})\bigr).
\]
\end{theorem}

\begin{proof}
\uses{thm:T-ad-one-ppow-eq, thm:T-sum-prime-mul-T-ad, lem:T-sum-prime, def:T-pp}
\leanok
Apply the key computation Theorem~\ref{thm:T-sum-prime-mul-T-ad} at exponent $k$ and rewrite the
two diagonal classes through the telescoping relation Theorem~\ref{thm:T-ad-one-ppow-eq}, using
$T(p,p^k)=T(p,p)\,T(1,p^{k-1})$ to express both $T(1,p^k)$, $T(1,p^{k+1})$ and $T(p,p^k)$ in
terms of the $T(p^\bullet)$. For $k=1$ the coefficient $c=p+1$ contributes an extra $+T(p,p)$
which combines with the base values $T(p^0)=1$ and $T(1,p)=T(p)$; for $k\ge 2$ one proceeds by
strong induction, substituting the recurrence already known at the smaller index and using that
$T(p)$ and $T(p,p)$ commute. Collecting the terms and isolating $T(p^{k+1})$ yields the displayed
three-term recurrence.
\end{proof}

\begin{theorem}[Shimura 3.24(3): prime-power product]\label{thm:T-sum-ppow-mul}
\lean{HeckeRing.GL2.T_sum_ppow_mul}
\uses{thm:T-sum-ppow-recurrence, def:T-pp, def:T-sum}
\leanok
Let $p$ be a prime and let $r\le s$. Then
\[
  T(p^r)\,T(p^s) \;=\; \sum_{i=0}^{r} p^i\,T(p^i,p^i)\,T(p^{r+s-2i}).
\]
\end{theorem}

\begin{proof}
\uses{thm:T-sum-ppow-recurrence, def:T-pp, def:T-sum}
\leanok
Strong induction on $r$. For $r=0$ the sum has the single term $T(p^s)$, and for $r=1$ the
statement is precisely the recurrence Theorem~\ref{thm:T-sum-ppow-recurrence} solved for
$T(p)\,T(p^s)$. For the inductive step write $T(p^{r+2})=T(p)\,T(p^{r+1})-p\,T(p,p)\,T(p^r)$ from
the recurrence and apply the induction hypothesis to $T(p^{r+1})\,T(p^s)$ and $T(p^r)\,T(p^s)$.
Each $T(p)\,T(p^{r+1+s-2i})$ splits, again by the recurrence, into a top term
$T(p^{r+2+s-2i})$ and a shifted term $p\,T(p,p)\,T(p^{r+s-2i})$; the shifted terms cancel against
the contribution of $-p\,T(p,p)\,T(p^r)$ except at the two endpoints, where reindexing
$T(p,p)^{i}\cdot T(p,p)=T(p^{i+1},p^{i+1})$ reassembles the sum over $i=0,\dots,r+2$. This is the
asserted product formula.
\end{proof}

\begin{theorem}[Coprime multiplicativity]\label{thm:T-sum-mul-coprime}
\lean{HeckeRing.GL2.T_sum_mul_coprime}
\uses{def:T-sum, def:T-ad, thm:commring-hecke-algebra}
\leanok
Let $m,n$ be positive integers with $(m,n)=1$. Then $T(m)\,T(n)=T(mn)$.
\end{theorem}

\begin{proof}
\uses{def:T-sum, def:T-ad, thm:commring-hecke-algebra}
\leanok
Expand both factors as sums over divisor pairs and multiply out:
$T(m)\,T(n)=\sum_{a\mid m}\sum_{b\mid n} T(a,m/a)\,T(b,n/b)$. Because $(m,n)=1$, the elementary
divisors of the two diagonal classes are coprime, so the corresponding diagonal double cosets
split as a single double coset under the Chinese Remainder Theorem, giving
$T(a,m/a)\,T(b,n/b)=T(ab,(m/a)(n/b))=T(ab,mn/(ab))$. The map $(a,b)\mapsto ab$ is a bijection from
divisor pairs of $(m,n)$ onto divisors of $mn$ (again by coprimality), so the double sum
reindexes to $\sum_{c\mid mn} T(c,mn/c)=T(mn)$. The terms where a divisibility constraint fails
contribute $0$ on both sides, matching up.
\end{proof}

\begin{theorem}[Shimura 3.24: general product]\label{thm:T-sum-mul}
\lean{HeckeRing.GL2.T_sum_mul}
\uses{thm:T-sum-ppow-mul, thm:T-sum-mul-coprime, def:T-pp, def:T-sum}
\leanok
For positive integers $m,n$,
\[
  T(m)\,T(n) \;=\; \sum_{d\,\mid\,(m,n)} d\,T(d,d)\,T\!\left(\frac{mn}{d^2}\right),
\]
the sum being over the positive divisors $d$ of the greatest common divisor $(m,n)$.
\end{theorem}

\begin{proof}
\uses{thm:T-sum-ppow-mul, thm:T-sum-mul-coprime, def:T-pp, def:T-sum}
\leanok
Both sides are multiplicative in the coprime factorisation of $m$ and $n$, so by induction on the
gcd it suffices to treat one prime at a time. Factor out a prime $q\mid(m,n)$, writing
$m=q^a m'$ and $n=q^b n'$ with $q\nmid m'n'$. The coprime multiplicativity
Theorem~\ref{thm:T-sum-mul-coprime} reduces $T(m)\,T(n)$ to $T(q^a)\,T(q^b)$ times the
$q$-free part $T(m')\,T(n')$, to which the induction hypothesis applies. The prime-power factor is
evaluated by Theorem~\ref{thm:T-sum-ppow-mul}, whose terms $q^i\,T(q^i,q^i)\,T(q^{a+b-2i})$ are
exactly the $d=q^i$ contributions to the divisor sum. Recombining the prime-power sum with the
$q$-free divisor sum, using that divisors of $(m,n)$ factor uniquely as $q^i$ times a divisor of
$(m',n')$ and that $T(d,d)$ is multiplicative on coprime $d$, reproduces the single sum over all
$d\mid(m,n)$.
\end{proof}

\begin{theorem}[Shimura 3.24(7): degree of $T(p^k)$]\label{thm:deg-T-sum-prime-pow}
\lean{HeckeRing.GL2.deg_T_sum_prime_pow}
\uses{def:deg, def:T-sum, def:T-ad}
\leanok
Let $p$ be a prime and $k\ge 0$. Then
\[
  \deg T(p^k) \;=\; \sum_{j=0}^{k} p^{j} \;=\; 1+p+p^2+\dots+p^k.
\]
\end{theorem}

\begin{proof}
\uses{def:deg, def:T-sum, def:T-ad}
\leanok
The degree homomorphism sends a diagonal class to its number of left cosets: for $2i<k$ one has
$\deg T(p^i,p^{k-i})=p^{k-2i-1}(p+1)$, while the scalar classes $T(p^i,p^i)$ have degree $1$.
Applying $\deg$ to the expansion $T(p^k)=\sum_{i} T(p^i,p^{k-i})$ and proceeding by strong
induction: for $k\ge 2$ split off the $i=0$ term, contributing $p^{k-1}(p+1)$, and observe that
each remaining term $T(p^{i+1},p^{k-1-i})$ has the same diagonal ratio — hence the same degree —
as the corresponding term of the $T(p^{k-2})$ expansion. Thus the tail equals
$\deg T(p^{k-2})=1+p+\dots+p^{k-2}$ by the induction hypothesis, and
$p^{k-1}(p+1)+(1+\dots+p^{k-2})=1+p+\dots+p^k$, the geometric sum.
\end{proof}

\begin{theorem}[Shimura 3.24(6): degree of $T(m)$]\label{thm:deg-T-sum}
\lean{HeckeRing.GL2.deg_T_sum}
\uses{def:deg, thm:deg-T-sum-prime-pow, thm:T-sum-mul-coprime, lem:deg-mul}
\leanok
For a positive integer $m$,
\[
  \deg T(m) \;=\; \sum_{d\,\mid\, m} d \;=\; \sigma_1(m),
\]
the sum of the positive divisors of $m$.
\end{theorem}

\begin{proof}
\uses{def:deg, thm:deg-T-sum-prime-pow, thm:T-sum-mul-coprime, lem:deg-mul}
\leanok
Both $\deg T(\cdot)$ and $\sigma_1$ are multiplicative arithmetic functions, so it suffices to
match them on prime powers and to use multiplicativity on coprime factors. On a prime power
Theorem~\ref{thm:deg-T-sum-prime-pow} gives $\deg T(p^k)=1+p+\dots+p^k=\sum_{d\mid p^k} d
=\sigma_1(p^k)$. For coprime $a,b$ the coprime multiplicativity $T(ab)=T(a)\,T(b)$
(Theorem~\ref{thm:T-sum-mul-coprime}) together with the ring-homomorphism property
$\deg(fg)=\deg f\,\deg g$ (Lemma~\ref{lem:deg-mul}) yields
$\deg T(ab)=\deg T(a)\,\deg T(b)=\sigma_1(a)\,\sigma_1(b)=\sigma_1(ab)$. Induction on the prime
factorisation of $m$ then gives $\deg T(m)=\sigma_1(m)$.
\end{proof}
